\documentclass[14pt,aspectratio=169,xcolor=dvipsnames]{beamer}
\usetheme{SimplePlus}
\usepackage{booktabs}
\usepackage{minted}

\title[short title]{Clase 15: Clases y polimorfismo I}
\subtitle{}
\author[NA Barnafi] {Nicolás Alejandro Barnafi Wittwer}
\institute[UC|CMM] 
{
    Pontificia Universidad Católica de Chile \\
    Centro de Modelamiento Matemático
}

\titlegraphic{
    \vspace{-1.8cm}
    \begin{flushright}
      \includegraphics[height=2.5cm]{../images/logos/puc.png} 
    \end{flushright}
}

\date{09/10/2024}
%\setbeamercovered{transparent}

\begin{document}
%%%%%%%%%%%%%%%%%%%%%%%%%%%%%%%%%%%%%%%%%%%%%%%%%%%%%%%
\begin{frame}
    \maketitle
\end{frame}
%%%%%%%%%%%%%%%%%%%%%%%%%%%%%%%%%%%%%%%%%%%%%%%%%%%%%%%
\begin{frame}[fragile]\frametitle{Motivación}
    \begin{minted}{python}
>>> df = pd.DataFrame() # Constructor
>>> type(df)
  <class 'pandas.core.frame.DataFrame'>
    \end{minted}
    \begin{itemize}
        \item Crear nuestros propios tipos de datos
        \item Crear abstracciones 
    \[ \text{animal $\in$ ser vivo, cuadrado $\in$ figura }\]
    \end{itemize}
\end{frame}
%%%%%%%%%%%%%%%%%%%%%%%%%%%%%%%%%%%%%%%%%%%%%%%%%%%%%%%
\begin{frame}[fragile]\frametitle{Clases conocidas}
    \begin{minted}{python}
>>> type(pd.DataFrame()) # pandas
    <class 'pandas.core.frame.DataFrame'>
>>> type(np.linspace(0,1)) # numpy
    <class 'numpy.ndarray'>
>>> type(2) # base
    <class 'int'>
    \end{minted}
\end{frame}
%%%%%%%%%%%%%%%%%%%%%%%%%%%%%%%%%%%%%%%%%%%%%%%%%%%%%%%
\begin{frame}[fragile]\frametitle{Sintaxis}
    \begin{minted}{python}
class A:
   # Método constructor
   # self es obligatorio
   def __init__(self, n):
      self.n = n # son distintos!
      self.n2 = n**2

a = A(4) # Constructor
print(a.n, a.n2)
print(type(a)) # <class '__main__.A'>
    \end{minted}

\end{frame}
%%%%%%%%%%%%%%%%%%%%%%%%%%%%%%%%%%%%%%%%%%%%%%%%%%%%%%%
\begin{frame}[fragile]\frametitle{Ejemplo: vector 2D}
    \begin{footnotesize}
    \begin{minted}{python}
class Vec2D:
   # Método constructor
   # self es obligatorio
   def __init__(self, x, y):
      self.x = x 
      self.y = y
   def __add__(self, v2): # magic function
      out = Vec2D(self.x, self.y) # copiar
      out.x += v2.x # agregar nuevos valores
      out.y += v2.y
      return out # retornar copia
   def __getitem__(self, i): # magic function
      if i == 0: return self.x
      if i == 1: return self.y
      else:
         print("ERROR") # i = 0 o 1
    \end{minted}
    \end{footnotesize}
\end{frame}
%%%%%%%%%%%%%%%%%%%%%%%%%%%%%%%%%%%%%%%%%%%%%%%%%%%%%%%
\begin{frame}\frametitle{Recap}
    \begin{itemize}
        \item Clases permiten crear abstracciones
        \item Todo objeto es una instancia de clase en Python
        \item Funciones mágicas permiten código legible!
    \end{itemize}
\end{frame}
%%%%%%%%%%%%%%%%%%%%%%%%%%%%%%%%%%%%%%%%%%%%%%%%%%%%%%%
\begin{frame}
    \maketitle
\end{frame}
%%%%%%%%%%%%%%%%%%%%%%%%%%%%%%%%%%%%%%%%%%%%%%%%%%%%%%%
\begin{frame}[noframenumbering]\frametitle{Mini ejercicios}
    \begin{itemize}
        \item Crear una clase Cuadrado que reciba como entrada el largo de sus aristas. Cree una función de clase que entregue el área del cuadrado.
        \item Extienda el ejemplo anterior al caso de un cubo.
    \begin{footnotesize}
    \begin{minted}{python}
class Cuadrado:
   def __init__(self, l):
   def area(self):
    \end{minted}
    \end{footnotesize}

    \end{itemize}
\end{frame}
%%%%%%%%%%%%%%%%%%%%%%%%%%%%%%%%%%%%%%%%%%%%%%%%%%%%%%%
\end{document}
